\clearpage
\SourceLeaf{411}%
\fancyhead[L]{\small VERSCHIEDENE AUFSÄTZE.}
\pdfbookmark[1]{Auflösung einer geometrischen Aufgabe.}{aufloesung-geometrische-1823}
\begin{center}
\rule{.8\textwidth}{.4pt}\\[4pt]
Astronomische Nachrichten. Nr. 42. 1823 November.\\[4pt]
\rule{.8\textwidth}{.4pt}\\[12pt]
\textit{Auflösung einer geometrischen Aufgabe.}
\end{center}
\hspace*{\parindent}\{In des Herrn Professors \textsc{Möbius} Beschreibung der Leipziger Sternwarte kommt S.\,61 eine kleine geometrische Aufgabe vor, nemlich:
\begin{quote}
Beliebige 5 Punkte \(A,B,C,D,E\) einer Ebene sind, je zwei, durch gerade Linien verbunden. Man kennt die somit entstehenden 5 Dreiecke \(EAB,ABC,BCD,CDE,DEA\) ihrem Inhalte nach, und verlangt daraus den Inhalt des Fünfecks \(ABCDE\).
\end{quote}
Herr Hofrath \textsc{Gauss}, der einige Wochen diesen Sommer bei mir verlebte, schrieb, wie er das Buch sah, folgende Auflösung hinein:

Man bezeichne die 5 Punkte mit \(1,2\ 3,4,5\), die
\begin{center}
\begin{tabular}{lrl}
Winkel&\(213\)&mit \(p\)\\
&\(214\)&— \(q\)\\
&\(215\)&— \(r\)\\\hline
Seiten&\(12\)&— \(t\)\\
&\(13\)&— \(u\)\\
&\(14\)&— \(v\)\\
&\(15\)&— \(w\)\\\hline
Dreiecke&\(123\)&— \(a\)\\
&\(234\)&— \(b\)\\
&\(345\)&— \(c\)\\
&\(451\)&— \(d\)\\
&\(512\)&— \(e\)\\
&\(124\)&— \(x\)\\
&\(134\)&— \(y\)\\
&\(135\)&— \(z\)\\\hline
Fünfeck&\(12345\)&— \(\omega\)
\end{tabular}
\end{center}
